% Homework URL:https://zhuoqing.blog.csdn.net/article/details/145879330

\section{Homework 2}

\subsection{直流和交流分解}

\task{必做题}

(1) $f(t)$是一个周期信号, 其直流分量$\displaystyle f_D(t)=\frac{1}{4}\int_{-2}^{2}f(t)\,\mathrm{d}t=-0.125$

(2) 同样是周期信号, 并且具有对称性.$\displaystyle f_D(t)=\frac{1}{\pi}\int_{0}^{\pi}\sin \frac{t}{2}\,\mathrm{d}t=\frac{1}{2\pi}$

(3) 可看作一个常数和一个积分为$0$的信号的叠加. $\displaystyle f_D(t)=\lim_{T\to\infty}(\frac{1}{2T}\int_{-T}^{T}f(t)\,\mathrm{d}t)=0.5$

(4) 周期信号且具有对称性. $\displaystyle f_D(t)=\frac{1}{0.5}\int_{0}^{0.5}4t^2\,\mathrm{d}t=\frac{1}{3}$.

\subsection{信号奇偶分解}

\task{奇偶分解}
如图\ref{fig:2-ess-1}所示.

\begin{figure*}[htbp]
    \centering
    \includegraphics[width=0.5\textwidth]{images/part2_essential.png}
    \caption{2.2.1波形}\label{fig:2-ess-1}
\end{figure*}

\task{求解原信号波形}
必做题如图\ref{fig:2.2.2-ess}所示.

\begin{figure*}[htbp]
    \centering
    \includegraphics[width=0.5\textwidth]{images/part2_solve.png}
    \caption{2.2.2波形-必做}\label{fig:2.2.2-ess}
\end{figure*}

\subsection{信号尺度变换}
必做题和选做题如图\ref{fig:2.3}.

\begin{figure*}[htbp]
    \centering
    \includegraphics[width=0.5\textwidth]{images/part3_transform.png}
    \caption{2.3波形}\label{fig:2.3}
\end{figure*}

\subsection{LTI系统响应}

\task{必做题}
当激励为$e_3(t)=\delta(t)=e_1'(t)$时, 由LTI系统的微分特性, 有:
\begin{align*}
    y_3(t)&=\frac{\mathrm d}{\mathrm dt}(\cos(3t)\cdot e^{-t}\cdot u(t)) \\
    &=\delta(t)-u(t)e^{-t}(\cos(3t)+3\sin(3t))
\end{align*}

激励$e_2=\delta(t-1)$为激励$e_3$的时变后信号, 因此
\begin{align*}
    y_2(t)=\delta(t-1)-u(t-1)e^{-t+1}(\cos(3(t-1)) + 3\sin(3(t-1)))
\end{align*}

\task{选做题}
(C)$x_2(t)=x_1(t)-x_1(t-2)$, 因此激励$y_2(t)=y_1(t)-y_1(t-2)$.

(D)$x_3(t)=x_1(t+1)+0.5x_1(t)$, 因此激励$y_3(t)=x_1(t+1)+0.5x_1(t)$.

波形如图\ref{fig:2.4-opt}.
\begin{figure*}[htbp]
    \centering
    \includegraphics[width=0.5\textwidth]{images/part4_LTI.png}
    \caption{2.4波形}\label{fig:2.4-opt}
\end{figure*}

\subsection{系统的可逆性}

\task{必做题}

(1)不可逆, $e_1[n]=1$和$e_2[n]=\cos(\pi t)$引起的输出相同, 均为$r[n]=1$.

(2)可逆, 逆系统为$r(t)=\displaystyle\frac{1}{2}\left.\frac{\mathrm d}{\mathrm d\tau}e(\tau)\right|_{\tau=\frac{t}{2}}$

(3)可逆, 逆系统为$r(t)=e(-t+2)$, 即系统本身.

(4)不可逆, $e_1(t)=2$和$e_2(t)=-2$均会给出输出$r(t)=8$.

\subsection{系统特性}

\begin{tabular}{|l|l|c|c|c|}
    \hline
    序号 & 系统输入输出关系 & 线性? & 时不变? & 因果? \\
    \hline 
    $1$ & $r(t)=e(2t)$ & 是 & 否 & 否 \\
    \hline 
    $2$ & $r(t)=\frac{\mathrm d}{\mathrm dt}e(t)$ & 是 & 是 & 否 \\
    \hline 
    $3$ & $r(t)=e(3-t)$ & 是 & 否 & 否 \\
    \hline 
    $4$ & $r(t)=e(0)$ & 是 & 否 & 否 \\
    \hline 
    $5$ & $r(t)=e(0.5t)$ & 是 & 否 & 否 \\
    \hline 
    $6$ & $r(t)=e(t)\cdot u(t)$ & 是 & 否 & 是 \\
    \hline 
    $7$ & $r(t)=t\cdot e(t)$ & 是 & 否 & 是 \\
    \hline
\end{tabular}

\subsection{从系统框图到方程}
(1)$x(t)=3y(t)+4y'(t)+y''(t)$.

(2)$x[n]=\frac{1}{2}y[n-1]+2y[n]+\frac{1}{2}y[n+1]$

(3)设第二个积分系统的输出信号为$f(t)$, 则:

\begin{equation*}
    \left\{
        \begin{array}{l}
            y(t)=f(t)-1.5f'(t) \\
            x(t) - f'(t) + 3f(t)=f''(t)
        \end{array}
    \right.
\end{equation*}

化简可得: $y''(t)+y'(t)+y(t)=-1.5x'(t)+x(t)$.

(4)设第二个延时系统的输出为$w[t]$, 则:

\begin{equation*}
    \left\{
        \begin{array}{l}
            w[t+2]=x[n]-w[t]-2w[t+1] \\
            y[n] = w[t]-3w[t+1]+1.5w[t+2]
        \end{array}
    \right.
\end{equation*}

化简可得: $1.5x[n+2]-3x[n+1]+x[n]=y[n+2]+2y[n+1]+y[n]$.
